\chapter{Conclusion, Limitations, and Future Work}
\label{chap:conclusion-limitations-and-future-work}

\section{Summary of the Dissertation}
\label{sec:8-1-summary-of-the-dissertation}

This dissertation examined whether the effectiveness of marketing content distributed through an agent-based WeChat network depends on the match between the content and the agents who share it. Using diffusion records from a single interior-design company on the WeiMai platform, the study linked LLM-assisted article-content features, corrected diffusion cascades, and observable agent attributes within a leakage-aware three-level design. The analysis covered 6,557 article–agent cases at the content level, 592 agents at the agent level, and 1,142 agent–topic pairs at the matching level.

Three findings followed. At the content level, title–content semantic consistency was the most robust correlate of diffusion and image richness the strongest of the continuous content predictors, but content alone explained only a modest share of the variation. At the agent level the differences were uneven: only Marketing/Operations agents stood out from the baseline, gender carried no consistent association once role and activity were accounted for, and the role that reached furthest also sat lowest in network centrality. At the matching level, which tested the central proposition most directly, stronger content–agent fit was associated with larger reach and fuller cascade shape, and the reach finding held under exposure weighting and the restriction to better-populated pairs. As at the agent level, however, better fit accompanied lower centrality, reproducing the divergence between reach-related and network-position dimensions.

Read together, the three levels support one conclusion: in this setting, similar content topics were associated with different diffusion outcomes depending on the agent context through which they were shared, so marketing effectiveness depended less on broadcasting the same content widely than on aligning content with the agents who carried it.

\section{Limitations}
\label{sec:8-2-limitations}

The most fundamental limitation is that the findings are associational, not causal. The analysis identified systematic relationships between content, agent, and matching variables and diffusion outcomes, but it did not manipulate any of these factors experimentally and cannot rule out that the observed associations are produced in part by unobserved factors. The study supports claims about association and provides a rationale for personalised assignment, but it does not establish that improving fit would cause diffusion to increase.

A second limitation is the single empirical context. All data came from one interior-design company on one SaaS platform, within the specific institutional and cultural setting of WeChat-based social marketing in China. The role categories, topic distribution, and audience behaviour observed here reflect that setting, and the magnitudes, and possibly the direction, of some associations may differ in other industries, platforms, or markets. The findings should be read as evidence from one well-documented case, not as established regularities, and their external validity remains an open question. Related to this, the study does not directly observe each agent's full audience network or the preferences of the recipients reached, so agent role is used as a proxy for audience context and not as a direct measure of it. The matching analysis consequently depends on the assumption, grounded in homophily, that an agent's professional role carries information about the audience the agent reaches.

A third limitation concerns the language-model-generated content features. Article topics and the title–content consistency measure were produced with the assistance of a large language model, which introduces classification and measurement uncertainty: topic assignments may be imperfect, and the consistency measure depends on the model's representation of meaning. Although the title–content consistency result was stable across alternative measures, the topic-based variables that underpin the matching analysis inherit whatever error the topic classification contains, so the matching findings rest partly on constructed features whose accuracy cannot be fully verified.

A fourth limitation concerns the construction of the matching measure itself. \varname{MatchScore} combines each article's topic profile with a set of job–topic relevance weights that the researcher specified from the functional interpretation of the agent role categories. These weights were neither estimated from the data nor validated against an external benchmark. Holding them independent of the diffusion outcomes protects the analysis against outcome leakage, but it also means the measure embeds a researcher judgement about which topics are relevant to which roles, so the Level 3 estimates, which carry the dissertation's central claim, are conditional on this weighting scheme. The binary \varname{ProfessionContentMatch} specification varies how fit is encoded but, being derived from the same weights, does not probe sensitivity to the weights themselves. Re-estimating the matching models under alternative weightings, or validating the weights against independent coders, would establish how far the central result depends on these choices.

A fifth limitation is the sparsity of the agent–topic cells at Level 3. Many agent–topic pairs were supported by only a small number of underlying cases. The reach association survived exposure weighting and the restriction to pairs with at least ten cases, but the stricter thirty-case sensitivity model could not be estimated with the full control specification, because the qualifying pairs were too concentrated within a single topic and across only two job categories. The matching results are better described as robust in sign and direction than as precisely estimated in magnitude.

A sixth limitation is the use of complete-case modelling. Cases with missing values on the modelled variables were excluded instead of imputed, which keeps the estimation transparent but reduces the effective sample and could introduce bias if the excluded cases differ systematically from those retained. The study did not attempt to recover these cases through imputation, so the estimates describe the complete-case sample and not the full set of recorded diffusion events.

A seventh limitation, methodological in nature, is that the reconstructed diffusion-derived variables cannot be interpreted as exogenous causal predictors. Measures such as structural virality, the Wiener index, repeat exposure, and centrality are computed from the same realised cascades that define the outcomes. The leakage-aware design responded to this by treating them as outcomes or descriptors, not as predictors. This protects the interpretation, but it also means the study cannot speak to whether these structural properties drive diffusion. They are used to characterise observed diffusion heterogeneity, not to explain it. Several of these measures, notably the Wiener index and reading duration at the agent level, were additionally weak or degenerate as outcomes and were retained only as diagnostics. In addition, the gender-assortativity measure is a custom bounded index, not the standard coefficient. It assigns zero both to cascades with balanced gender mixing and to those with no usable gendered ties, so its agent-level averages are best read as approximate descriptors of same-gender tie preference and not as precise homophily estimates.

Finally, the data record diffusion behaviours, not downstream business outcomes such as inquiries, purchases, or revenue. The findings speak to marketing diffusion and not to sales conversion.

\section{Future Research}
\label{sec:8-3-future-research}

The most direct extension concerns causal identification. Because the present findings are associational, the natural next step is to test the matching proposition experimentally. A field experiment that randomly assigned comparable articles, or the same article where feasible, to better- or worse-matched agents would separate the effect of fit from the unobserved conditions that observational data cannot exclude. It would establish whether improving content–agent fit causes diffusion to increase or merely accompanies it. Where full randomisation is impractical, quasi-experimental designs that exploit naturally occurring variation in assignment could provide partial identification. Such studies would convert the personalised-assignment rationale developed here from an evidence-informed expectation into a tested intervention.

Since the data came from a single interior-design firm on a single platform, extending the analysis to other industries, other SaaS platforms, and other markets would test whether the matching relationship, and the divergence between reach-related and network-position dimensions, generalise beyond this case. Comparative work across firms whose agents differ in role structure and audience composition would also clarify which features of the present setting are essential to the findings and which are incidental to it.

A third direction follows from the measurement limitations. Future studies could observe audience context more directly, for example through recipient-level attributes or audience-network information where these are available, instead of relying on agent role as a proxy. They could also strengthen the content features by validating language-model-generated topics and consistency scores against human coding or alternative models, and by quantifying how measurement error in these features propagates into the matching estimates. They could likewise probe the job–topic relevance weights that enter \varname{MatchScore}, re-estimating the matching models under alternative weightings or validating the weights against independent coders, since these weights are the second input to the fit measure alongside the language-model topics. Improved measurement on either side of the fit relationship would sharpen the matching analysis and reduce its dependence on constructed proxies.

The present study observed diffusion behaviours, not downstream business results. Linking diffusion measures to inquiries, purchases, or revenue would establish whether better content–agent fit ultimately improves commercial outcomes and not only reach. Richer data would also support better-populated agent–topic cells, stricter sparse-cell sensitivity checks, and missing-data strategies beyond complete-case exclusion.

Taken together, these directions describe a programme that would move from the associational, single-context, diffusion-focused evidence presented here towards causal, multi-context, and outcome-linked tests of the same central proposition. The contribution of the present dissertation is to establish that proposition as empirically credible in a real agent-based marketing setting, and to provide the framework, measures, and findings on which such further work can build.
